\subsubsection{When $\Dcontr$ has a long path}\label{sec:biclique:longpath}

In this section, we prove Lemma~\ref{lem:digging} when the length of a longest path in $\Dcontr$ is at least $2t+2$.
Let $P=(v_1, \ldots,v_p)$
be a directed longest path in $\Dcontr$ (then $p \geq 2t+2$).
Recall that none of the edges of $\Dcontr$, and in particular, none  of the edges of $P$, are {\em contraction-redundant} with respect to $(A,B)$.

\begin{lemma}\label{claim:biclique1}
For each $e \in E(P)$, $\Dcontr/e$ is acyclic.
\end{lemma}

\begin{proof}
Suppose, for the sake of contradiction, that there is a cycle in $\Dcontr/e$. Let $e=(v_i,v_{i+1})$. 
Then, either there is a $(v_i,v_{i+1})$-path in $\Dcontr - \{(v_i,v_{i+1})\}$ or a $(v_{i+1},v_i)$-path in $\Dcontr$. 
If there exists a $(v_{i+1},v_i)$-path in $\Dcontr$ then together with the edge $e=(v_i,v_{i+1})$, it creates a directed closed walk in $\Dcontr$. 
This contradicts that $\Dcontr$ is acyclic. 

In the other case, suppose there exists a $(v_i,v_{i+1})$-path, say $P'$ in $\Dcontr -\{(v_i,v_{i+1})\}$. We first claim that the internal vertices of $P'$ are disjoint from $P$. Suppose not, let $v_j \in V(P)$ be the first internal vertex of $P'$ that belongs to $P$. Then there is a $(v_i,v_j)$-path in $\Dcontr$ and $(v_j,v_{i+1})$-path in $\Dcontr$. If $j<i$, then the $(v_i,v_j)$-path together with the $(v_j,v_i)$-subpath of $P$, gives a closed walk in $\Dcontr$, contradicting its acyclicity. Otherwise, $j>i+1$. In this case, the $(v_j,v_{i+1})$-path together with the $(v_{i+1},v_j)$-subpath of $P$ gives a closed walk in $\Dcontr$, again contradicting its acyclicity.
\end{proof}

Since each edge of $\Dcontr$, in particular, each edge of $P$, is not contraction-redundant, from Lemma~\ref{claim:biclique1} and the definition of contraction-redundant, for each $e \in E(P)$ either
\begin{enumerate}
\item there exists an $aa$-path, or a $bb$-path in $\Dcontr/e$ (that is, $(\Dcontr/e)^{\star}$ does not contain the $(A,B)$ -induced-biclique), or 
\item there exists $i \in [g(t)]$ an $a_{i}b_{i}$-path in $(\Dcontr/e )- \{(a_{i},b_{i})\}$ ($(a_{i},b_{i})$ is not a minimal edge in $\Dcontr/e$).
\end{enumerate}

\begin{lemma}\label{obs:biclique1}
Let $e=(v_i,v_{i+1}) \in E(P)$ such that $H_{contr}/e$ has an $(a_j,b_j)$-path. Then there exists an $(a_j,v_{i+1})$-path and a $(v_i,b_j)$-path in $\Dcontr$.

If $\Dcontr/e$ has a $(a_j,a_{\ell})$-path, for some $a_j,a_{\ell} \in A$, then there exists an $(a_j, v_{i+1})$-path and a $(v_i,a_{\ell})$-path in $\Dcontr$.

If $\Dcontr/e$ has a $(b_j,b_{\ell})$-path, for some $b_j,b_{\ell} \in B$, then there exists an $(b_j, v_{i+1})$-path and $(v_i,b_{\ell})$-path in $\Dcontr$.
\end{lemma}

\begin{proof}
Let $e=(v_i,v_{i+1}) \in E(P)$ such that $D_{contr}/e$ has an $(a_j,b_j)$-path. Then either there exists an $(a_j,v_i)$-path and a $(v_{i+1},b_j)$-path in $\Dcontr$, or there exists an $(a_j,v_{i+1})$-path and $(v_i,b_j)$-path in $\Dcontr$. In the later case, we are done. The former case implies an $(a_j,b_j)$-path in $\Dcontr$ of length at least three (with at least three edges), which contradicts that $(a_j,b_j)$ is a minimal edge in $\Dcontr$.

If $D_{contr}/e$ has an $(a_j,a_{\ell})$-path, then either there exists an $(a_j,v_i)$-path and a $(v_{i+1},a_{\ell})$-path in $\Dcontr$, or there exists an $(a_j,v_{i+1})$-path and $(v_i,a_{\ell})$-path in $\Dcontr$. In the later case, we are done. The former case implies an $(a_j,a_{\ell})$-path in $\Dcontr$ which contradicts that $(A,B)$ is an induced biclique in $\Dcontr^{\star}$ (in particular, that $A$ is an independent set in $\Dcontr^{\star}$).
One can similarly show that if $D_{contr}/e$ has a $(b_j,b_{\ell})$-path, then there exists an $(b_j,v_{i+1})$-path and $(v_i,b_{\ell})$-path in $\Dcontr$.
\end{proof}

For each $e=(v_i,v_{i+1}) \in E(P)$, we say that:
\begin{enumerate}
\item $e$ is {\em $j$-irredundant} in $\Dcontr$, if there is an $(a_j,v_{i+1})$-path and a $(v_i,b_j)$-path in $\Dcontr$.
\item $e$ is {\em $aa$-irredundant} in $\Dcontr$, if there exists $a_j,a_{\ell} \in A$ such that there is an $(a_j,v_{i+1})$-path and a $(v_i,a_{\ell})$-path in $\Dcontr$.
\item  $e$ is {\em $bb$-irredundant} in $\Dcontr$, if there exists $b_j,b_{\ell} \in B$ such that there is an $(b_j,v_{i+1})$-path and $(v_i,b_{\ell})$-path in $\Dcontr$.
\end{enumerate}

From Lemmas~\ref{claim:biclique1} and~\ref{obs:biclique1}, for each $e \in E(P)$, either $e$ is $j$-irredundant, for some $j \in [g(t)]$, or is $aa$-irredundant or is $bb$-irredundant.

\begin{lemma}\label{claim:biclique21}
Let $e=(v_i,v_{i+1}) \in E(P)$ such that $e$ is $aa$-irredundant in $\Dcontr$. Then $i=1$.
\end{lemma}

\begin{proof}
For the sake of contradiction, suppose that $i>1$. Let $e'=(v_{i-1},v_i)$. 
If $e'$ is $aa$-irredundant or $e'$ is $j$-irredundant for some $j \in [g(t)]$, then there exists a path from some vertex of $A$ to $v_i$ in $\Dcontr$. Also since $e$ is $aa$-irredundant, there exists a path from $v_i$ to some vertex of $A$ in $\Dcontr$. This implies either a directed closed walk in $\Dcontr$, contradicting its acyclicity, or a path between two distinct vertices of $A$ in $\Dcontr$, contradicting that $A$ is an independent set in $\Dcontr^{\star}$.

If $e'$ is $bb$-irredundant, then there exists a path from some vertex of $B$ to $v_i$. Also since $e$ is $aa$-irredundant, there exists a path from $v_i$ to some vertex of $A$. This implies a path from some vertex of $B$ to some vertex of $A$. Since $(A,B)$ is an induced biclique in $\Dcontr^{\star}$, that is there is a path from every vertex of $A$ to every vertex of $B$ in $\Dcontr$, this implies a directed closed walk in $\Dcontr$, contradicting its acyclicity.
\end{proof}

\begin{lemma}\label{claim:biclique22}
Let $e=(v_i,v_{i+1}) \in E(P)$ such that $e$ is $bb$-irredundant in $\Dcontr$. Then $i+1=p$.
\end{lemma}

\begin{proof}
For the sake of contradiction, suppose that $i+1<p$. Let $e'=(v_{i+1},v_{i+2})$. 
If $e'$ is $bb$-irredundant or $e'$ is $j$-irredundant for some $j \in [g(t)]$, then there exists a path from $v_{i+1}$ to some vertex of $B$. Also since $e$ is $bb$-irredundant, there exists a path from some vertex of $B$ to $v_{i+1}$. This either implies a directed closed walk in $\Dcontr$, contradicting its acyclicity, or a path between two distinct vertices of $B$ in $\Dcontr$, contradicting that $B$ is an independent set in $\Dcontr^{\star}$.

If $e'$ is $aa$-irredundant, then there exists a path from $v_{i+1}$ to some vertex of $A$ in $\Dcontr$. Also since $e$ is $bb$-irredundant, there exists a path from some vertex of $B$ to $v_{i+1}$ in $\Dcontr$. This implies a path from some vertex of $B$ to some vertex of $A$ in $\Dcontr$. Since 
$(A,B)$ is an induced-biclique in $\Dcontr^{\star}$, this implies a directed closed walk
in $\Dcontr$, contradicting its acyclicity.
\end{proof}

Let $P'$ be a subpath obtained from $P$ by without the first and last edge of $P$, that is $P'=(v_2, \ldots,v_{p-1})$. From Lemmas~\ref{claim:biclique21} and~\ref{claim:biclique22}, each edge $e \in E(P')$ is $j$-irredundant for some $j \in [g(t)]$.

\begin{lemma}\label{claim:biclique3}
Let $e,e'$ be distinct edges of $P'$ such that $e$ is $q$-irredundant in $\Dcontr$ and $e'$ is $r$-irredundant in $\Dcontr$, for some $q,r \in [g(t)]$. Then $q \neq r$.
\end{lemma}

\begin{proof}
For the sake of contradiction, suppose that $q=r$. Without loss of generality, let $e=(v_i,v_{i+1})$ and $e'=(v_j,v_{j+1})$ such that $i<j$. Then there exists an $(a_q,v_{i+1})$-path and $(v_j,b_q)$-path in $\Dcontr$. Since $i < j$, there exists a path from $(v_i,v_j)$ (which is the the $(v_{i+1},v_j)$-subpath of $P'$) in $\Dcontr$. This implies an $(a_q,b_q)$-path of length strictly greater than one in $\Dcontr$, contradicting that $(a_q,b_q)$ is a minimal edge in $\Dcontr$.
\end{proof}

For each $e=(v_i,v_{i+1}) \in E(P')$, fix a $\rho(i) \in [g(t)]$, such $e$ is $\rho(i)$-irredundant in $\Dcontr$. 
The following observation follows from Lemma~\ref{claim:biclique3}.

\begin{observation}\label{obs:biclique1}
For each $v_i \in V(P')$, $i \in [3, p-2]$, there exists a $(a_{\rho(i-1)},v_i)$-path and a $(v_i, b_{\rho(i)})$-path in $\Dcontr$.  Also, for $i,j \in [3,p-2]$, $i \neq j$, $\rho(i) \neq \rho(j)$.
\end{observation}

For each $i \in [3,p-2]$, let $P^{\aaa}_i$ denote an (arbitrarily) fixed $(a_{\rho(i-1)},v_i)$-path in $\Dcontr$ and let $P^{\bbb}_i$ denote an (arbitrarily) fixed $(v_i,b_{\rho(i)})$-path in $\Dcontr$.
Let $E_1 = (v_3, b_{\rho(3)}), \allowbreak  (v_4, b_{\rho(4)}), \ldots, \allowbreak (v_{p/2-2}, b_{\rho(p/2-2)})$ and let 
 $E_2 = ((a_{\rho(p/2-2)},v_{p/2-1}), \allowbreak  (a_{\rho(p/2-1)}, \allowbreak v_{p/2}), \ldots, (a_{\rho(p-1)}, \allowbreak v_{p-2})$. 
 From Observation~\ref{obs:biclique1}, each edge of $E_1 \cup E_2$ is an edge in $D^{\star}_{\contr}$ and the endpoints of the edges in $E_1 \cup E_2$ are distinct. We will now show that $(E_1,E_2)$ is an ordered $(p/2-1)$-\toughpair\ in $D^{\star}_{\contr}$. Since $p \geq 2t+2$, this proves Lemma~\ref{lem:digging} in the case when the length of a longest path in $\Dcontr$ is at least $2t+2$.
 
\begin{lemma}\label{lem:long-path-ordered-tough-pair}
$(E_1,E_2)$ is an ordered $(p/2-1)$-\toughpair\ in $D^{\star}_{\contr}$.
\end{lemma}

\begin{proof}
To prove the lemma, we prove the following claims.

\begin{claim}[Weak independence of $E_1$ (resp.~$E_2$)]\label{claim:biclique-long-path-weak}
The edges in $E_1$ are pairwise weakly independent. Similarly, the edges in $E_2$ are pairwise weakly independent in $D^{\star}_{\contr}$.
\end{claim}

\begin{claimproof}
Let $(a_{\rho(i-1)},{v}_{i}), \allowbreak (a_{\rho(j-1)},{v}_{j}) \in E_1$. For the sake of contradiction, say $(v_j, a_{\rho(i-1)}) \in \Dcontr^{\star}$. Then this implies a $(a_{\rho(j-1},a_{\rho(i-1)})$-path in $\Dcontr^{\star}$ contradicting that $A$ is an independent set in $\Dcontr^{\star}$.
Using symmetric arguments one can show that all edges in $E_2$ are pairwise weakly independent.
\end{claimproof}

\begin{claim}[Orderedness on $E_1$ and $E_2$]\label{claim:biclique-long-path-ordered}
For any $i,j \in [3, p-2]$, $i <j$, $(a_{\rho(j-1)}, v_{i}), \allowbreak (a_{\rho(j-1)}, a_{\rho(i-1)}),  \allowbreak ({v}_{j},  \allowbreak b_{\rho(i)}), \allowbreak (v_j,v_i) \not \in E(\Dcontr^{\star})$.
\end{claim}

\begin{claimproof}
Since $A$ is an independent set in $\Dcontr^{\star}$, $(a_{\rho(j-1)}, a_{\rho(i-1)}) \not \in E(\Dcontr^{\star})$. Since $\Dcontr^{\star}$ is acyclic and there is a $(v_i,v_j)$-path in $\Dcontr$, we conclude that $(v_j,v_i) \not \in E(\Dcontr^{\star})$.
If there is an $(a_{\rho(j-1)}, {v}_{i})$-path in $\Dcontr$, then this together with the $( {v}_{i},v_{j-1})$-subpath of $P'$ and the $(v_{i-1},b_{\rho(i-1)})$ path $P^{\bbb}_{i-1}$, implies a $(a_{\rho(j-1),b_{\rho(j-1)}})$-path of length strictly greater than $1$ in $\Dcontr$, which is a contradiction. Using symmetric arguments, one can show that 
$({v}_{j}, b_{\rho(i)}), (v_j,v_i) \not \in E(\Dcontr^{\star})$.
\end{claimproof}

\begin{claim}[Weak independence between $E_1$ and $E_2$]\label{claim:biclique-long-path-weak-between}
Let $(v_i,b_{\rho(i)}) \in E_1$ and $(a_{\rho(j-1)},v_j) \in E_2$.
There is $(v_j,b_{\rho(i)}), (b_{\rho(i)},a_{\rho(i-1)}) \not \in E(\Dcontr^{\star})$. 
\end{claim}

\begin{claimproof}
First observe from the construction of $E_1,E_2$ that $i <j$.
If $(v_j,b_{\rho(i)}) \in E(\Dcontr^{\star})$, then the $(a_{\rho(i)},v_{i+1})$-path $P^{\aaa}_{i+1}$, together with the $(v_{i+1},v_j)$-subpath of $P'$, together with the $(v_j,b_{\rho(i)})$-path, implies a $(a_{\rho(i)},b_{\rho(i)})$-path of length strictly greater than $1$ in $\Dcontr$, which is a contradiction. 

If $(b_{\rho(i)},a_{\rho(i-1)})  \in E(\Dcontr^{\star})$, then the $(a_{rho(i-1)},b_{\rho(i)})$-path obtained by concatenating $P^{\aaa}_i$ and $P^{\bbb}_i$, together with the $(b_{\rho(i)},a_{\rho(i-1)})$-path, implies a $(a_{\rho(i-1)},a_{\rho(j-1)})$-path in $\Dcontr$, which contradicts that $A$ is an independent set in $\Dcontr$.
\end{claimproof}

\begin{claim}[Strong independence between $E_1$ and $E_2$]\label{claim:biclique-long-path-strong}
Let $(v_i,b_{\rho(i)}) \in E_1$ and $(a_{\rho(j-1)},v_j) \in E_2$.
Then $(v_i,a_{\rho(j-1}), \allowbreak (a_{\rho(j-1)},v_i), \allowbreak (b_{\rho(i)},v_j), \allowbreak (v_j,b_{\rho(i)}) \not \in E(\Dcontr^{\star})$.
\end{claim}

\begin{claimproof}
Recall that $i >j$.
If $(v_i,a_{\rho(j-1)}) \in E(\Dcontr^{\star})$, then together with $P^{\aaa}_i$, it contradicts that $A$ is independent in $\Dcontr^{\star}$. Similarly, if $(b_{\rho(i)},v_j) \in E(\Dcontr^{\star})$, then it contradicts the independence of $B$ in $\Dcontr^{\star}$.

If $(a_{\rho(j-1)},v_i) \in E(\Dcontr^{\star})$, then this together with the $(v_i,v_{j-1})$-subpath of $P'$, together with the path $P^{\bbb}_{j-1}$, implies a $(a_{\rho(j-1)},b_{\rho(j-1)})$-path of length strictly greater than $1$ in $\Dcontr$, which is a contradiction. Similarly if, $(v_j,b_{\rho(i)}) \in E(\Dcontr^{\star})$, then this implies a $(a_{\rho(i)},b_{\rho(i)})$-path of length strictly greater than $1$ in $\Dcontr$, which is a contradiction.
\end{claimproof}
This concludes the proof.
\end{proof}